\documentclass[12pt,a4paper]{article}
\newcommand{\LabNumber}{5}
\newcommand{\LabTitle}{Criptografia cu chei publice:\\ RSA, ElGamal și Diffie–Hellman}
\newcommand{\LabTitleShort}{Lucrarea de laborator nr. 5 — Criptografie \& Securitate Informațională}
% ---------------------------------------------------------------------------
% Common preamble for the CS laboratory reports (compile with XeLaTeX)
% Strunga Daniel-Ioan, FAF-242 — https://github.com/don09dandidi/CS
% ---------------------------------------------------------------------------
\usepackage{fontspec}
\setmainfont{Latin Modern Roman}
\setsansfont{Latin Modern Sans}
\setmonofont{DejaVu Sans Mono}[Scale=0.85]
\usepackage[english]{babel}
\babelprovide[import, main]{romanian}
\usepackage[a4paper, margin=2.5cm, headheight=15pt]{geometry}
\usepackage{amsmath, amssymb}
\usepackage{booktabs, array, longtable, tabularx}
\usepackage{xcolor}
\usepackage{listings}
\usepackage{fvextra}
\usepackage[most]{tcolorbox}
\usepackage{fancyhdr}
\usepackage{tikz}
\usetikzlibrary{arrows.meta, positioning, shapes.geometric, fit, calc}
\usepackage{enumitem}
\usepackage{caption}
\usepackage{float}
\usepackage[hidelinks]{hyperref}
\usepackage{url}

\definecolor{codebg}{RGB}{248,248,250}
\definecolor{codeframe}{RGB}{200,200,210}
\definecolor{kw}{RGB}{0,0,160}
\definecolor{cmt}{RGB}{0,120,0}
\definecolor{str}{RGB}{160,30,30}
\definecolor{termbg}{RGB}{30,32,40}
\definecolor{termfg}{RGB}{225,228,235}
\definecolor{termbar}{RGB}{60,63,75}
\definecolor{boxblue}{RGB}{230,238,250}

\lstdefinestyle{py}{
  language=Python, basicstyle=\ttfamily\small, backgroundcolor=\color{codebg},
  frame=single, rulecolor=\color{codeframe}, numbers=left, numberstyle=\tiny\color{gray},
  keywordstyle=\color{kw}\bfseries, commentstyle=\color{cmt}\itshape, stringstyle=\color{str},
  breaklines=true, showstringspaces=false, columns=fullflexible, keepspaces=true, tabsize=4,
  extendedchars=true, inputencoding=utf8
}
\lstdefinestyle{sh}{
  language=bash, basicstyle=\ttfamily\small, backgroundcolor=\color{codebg},
  frame=single, rulecolor=\color{codeframe}, numbers=left, numberstyle=\tiny\color{gray},
  keywordstyle=\color{kw}\bfseries, commentstyle=\color{cmt}\itshape, stringstyle=\color{str},
  breaklines=true, showstringspaces=false, columns=fullflexible, keepspaces=true
}
\lstdefinestyle{wl}{
  basicstyle=\ttfamily\small, backgroundcolor=\color{codebg},
  frame=single, rulecolor=\color{codeframe}, numbers=left, numberstyle=\tiny\color{gray},
  morekeywords={RandomPrime,RandomInteger,PowerMod,Mod,GCD,FromDigits,IntegerDigits,
    ToCharacterCode,FromCharacterCode,IntegerLength,BitLength,Module,Print,Hash,Length},
  keywordstyle=\color{kw}\bfseries, commentstyle=\color{cmt}\itshape,
  morecomment=[s]{(*}{*)}, breaklines=true, columns=fullflexible, keepspaces=true
}

% Console session box (as in the Lab 2 report)
\newtcolorbox{terminal}[1][Sesiune Terminal Consolă]{
  enhanced, breakable, colback=termbg, colframe=termbar, coltext=termfg,
  fonttitle=\small\ttfamily\bfseries, title={\textbullet\textbullet\textbullet\quad #1},
  left=4pt, right=4pt, top=2pt, bottom=2pt, boxrule=0.6pt, arc=2pt
}
\DefineVerbatimEnvironment{termout}{Verbatim}{fontsize=\scriptsize, breaklines=true, breakanywhere=true}

\newtcolorbox{infobox}[1]{
  enhanced, breakable, colback=boxblue, colframe=kw!60, fonttitle=\bfseries,
  title={#1}, boxrule=0.6pt, arc=2pt
}

\newcommand{\bignum}[1]{\par\nobreak\vspace{1pt}{\noindent\ttfamily\footnotesize\seqsplit{#1}\par}\vspace{5pt}}
\usepackage{seqsplit}

\pagestyle{fancy}
\fancyhf{}
\renewcommand{\headrulewidth}{0.4pt}
\fancyhead[L]{\footnotesize\textit{\LabTitleShort}}
\fancyhead[R]{\footnotesize strunga.daniel-ioan}
\fancyfoot[C]{\thepage}

\newcommand{\makereporttitle}{%
\begin{titlepage}
  \centering
  {\large\bfseries UNIVERSITATEA TEHNICĂ A MOLDOVEI\par}
  \vspace{2pt}
  {\normalsize Facultatea Calculatoare, Informatică și Microelectronică\par}
  {\normalsize Departamentul Ingineria Software și Automatică\par}
  \vspace{5cm}
  {\LARGE\bfseries LUCRAREA DE LABORATOR NR. \LabNumber\par}
  \vspace{8pt}
  {\Large\bfseries \LabTitle\par}
  \vspace{10pt}
  {\normalsize Disciplina: Criptografie și Securitate Informațională\par}
  \vfill
  \begin{tabular*}{\textwidth}{@{}l@{\extracolsep{\fill}}r@{}}
    \textbf{Autor:} & \textbf{Verificat de:} \\
    Daniel-Ioan Strunga & Lector / Supervizor \\
    Student gr. FAF-242 & Maia Zaica \\
  \end{tabular*}
  \vspace{3cm}

  {\normalsize Chișinău — 2026\par}
\end{titlepage}
}


\begin{document}
\makereporttitle
\tableofcontents
\newpage

\section{Introduction and Laboratory Objectives}
Public-key (asymmetric) cryptography, introduced by Diffie and Hellman in 1976, removes the need
for a secret channel to distribute keys: every participant owns a key pair — a public key that
can be published and a private key that is never disclosed. Its security rests on
computationally hard number-theoretic problems: \textbf{integer factorisation} (RSA) and the
\textbf{discrete logarithm problem} in $\mathbb{Z}_p^*$ (ElGamal, Diffie–Hellman).

Objectives of the laboratory work:
\begin{enumerate}
  \item \textbf{Task 2.1} — generate RSA keys with a modulus of at least 2048 bits, encrypt and
        decrypt the message $m=$ \textit{„Nume Prenume”} $=$ \texttt{"<<message>>"}.
  \item \textbf{Task 2.2} — generate ElGamal keys for the given 2048-bit prime $p$ and generator
        $g=2$, encrypt and decrypt the same message.
  \item \textbf{Task 3} — perform the Diffie–Hellman key exchange between Alice and Bob, who will
        use AES with a 256-bit key; the secrets $a$ and $b$ are chosen randomly according to the
        requirements of the algorithm.
  \item Comment in detail every step of the algorithms.
\end{enumerate}

\begin{infobox}{Tooling note}
The laboratory guide suggests Wolfram Alpha / Mathematica. All computations in this report were
performed with the script \texttt{lab5/public\_key\_crypto.py}, which uses only Python's
arbitrary-precision integers and the exact equivalents of the Wolfram functions
(\texttt{RandomPrime}, \texttt{RandomInteger}, \texttt{PowerMod}, \texttt{Mod}). The same steps
written in Wolfram Language are given in \texttt{lab5/wolfram/lab5.wl} and in
Appendix~\ref{app:wl}; random values will of course differ from run to run.
\end{infobox}

\section{Message Encoding}
As required, the message is converted to a decimal number through the hexadecimal representation
of its ASCII codes:

\begin{table}[H]
\centering
\caption{ASCII encoding of the message}
\small
\setlength{\tabcolsep}{3pt}
\begin{tabular}{l*{14}{c}}
\toprule
<<ascii_table>>
\bottomrule
\end{tabular}
\end{table}

\begin{align*}
  \text{hex}(m) &= \texttt{<<hex>>}\\
  m &= <<m>> \quad (<<m_bits>>\text{ bits})
\end{align*}
Since $m < n$ (RSA) and $m < p$ (ElGamal), the message is encrypted as a single block. Decoding is
the reverse process: the decimal number is written in base 256 and every digit is an ASCII code.

\section{Task 2.1 — RSA}
\subsection{Algorithm}
\begin{enumerate}
  \item Choose two large random primes $p$ and $q$ of 1024 bits each (the two most significant
        bits set, so that $n$ has exactly 2048 bits).
  \item Compute the modulus $n = p\cdot q$ and Euler's function $\varphi(n)=(p-1)(q-1)$.
  \item Choose the public exponent $e$ with $1<e<\varphi(n)$, $\gcd(e,\varphi(n))=1$; the
        standard value $e = 2^{16}+1 = 65537$ is used.
  \item Compute the private exponent $d = e^{-1} \bmod \varphi(n)$ (extended Euclid).
  \item Public key: $(e,n)$; private key: $(d,n)$.
  \item Encryption: $c = m^{e} \bmod n$; decryption: $m = c^{d} \bmod n$.
\end{enumerate}
Correctness follows from Euler's theorem: $c^d = m^{ed} = m^{1+t\varphi(n)} \equiv m \pmod n$.

\subsection{Key Generation}
\noindent\textbf{Prime} $p$ (1024 bits):
\bignum{<<rsa.p>>}
\noindent\textbf{Prime} $q$ (1024 bits):
\bignum{<<rsa.q>>}
\noindent\textbf{Modulus} $n=p\cdot q$ (\textbf{<<rsa.n_bits>> bits}):
\bignum{<<rsa.n>>}
\noindent\textbf{Euler's function} $\varphi(n)=(p-1)(q-1)$:
\bignum{<<rsa.phi>>}
\noindent\textbf{Public exponent:} $e = <<rsa.e>>$, \quad $\gcd(e,\varphi(n)) = 1$.\\[4pt]
\noindent\textbf{Private exponent} $d = e^{-1} \bmod \varphi(n)$:
\bignum{<<rsa.d>>}
Check: $e\cdot d \bmod \varphi(n) = <<rsa.check_ed>>$.

\subsection{Encryption and Decryption}
\noindent\textbf{Ciphertext} $c = m^{e} \bmod n$:
\bignum{<<rsa.c>>}
\noindent\textbf{Decryption} $c^{d} \bmod n$:
\bignum{<<rsa.m_decrypted>>}
which is equal to $m$; converting it back to ASCII gives \texttt{"<<message>>"}.

\section{Task 2.2 — ElGamal}
\subsection{Algorithm}
The public parameters are the given prime $p$ (2048 bits) and the generator $g=2$. The prime is a
\emph{safe prime} ($p=2q'+1$ with $q'$ prime — verified with the Miller–Rabin test), so the
discrete logarithm problem in $\mathbb{Z}_p^*$ has no small subgroups to exploit.
\begin{enumerate}
  \item \textbf{Key generation (receiver):} choose the private key $x$ randomly,
        $1 < x < p-1$, and publish $y = g^{x} \bmod p$. Public key: $(p,g,y)$.
  \item \textbf{Encryption (sender):} choose a fresh random ephemeral key $k$, $1<k<p-1$, and
        compute
        $c_1 = g^{k} \bmod p$, \; $c_2 = m\cdot y^{k} \bmod p$. The ciphertext is $(c_1,c_2)$.
  \item \textbf{Decryption (receiver):} the shared value is $s=c_1^{x}=g^{kx}$, and
        $m = c_2 \cdot s^{-1} \bmod p = c_2 \cdot c_1^{\,p-1-x} \bmod p$ (Fermat: $s^{-1}=s^{p-2}$,
        and $c_1^{p-1-x}=c_1^{-x}$).
\end{enumerate}
The ephemeral key $k$ must never be reused: from two ciphertexts with the same $k$ one obtains
$m_1/m_2 = c_2/c_2'$.

\subsection{Values}
\noindent\textbf{Prime} $p$ (given, 2048 bits) and $g = 2$:
\bignum{<<P>>}
\noindent\textbf{Private key} $x$:
\bignum{<<eg.x>>}
\noindent\textbf{Public key} $y = g^{x} \bmod p$:
\bignum{<<eg.y>>}
\noindent\textbf{Ephemeral key} $k$:
\bignum{<<eg.k>>}
\noindent\textbf{Ciphertext} $c_1 = g^{k} \bmod p$:
\bignum{<<eg.c1>>}
\noindent\textbf{Ciphertext} $c_2 = m \cdot y^{k} \bmod p$:
\bignum{<<eg.c2>>}
\noindent\textbf{Decryption} $m = c_2 \cdot c_1^{\,p-1-x} \bmod p$:
\bignum{<<eg.m_decrypted>>}
The decrypted number equals $m$ and decodes to \texttt{"<<message>>"}.

\section{Task 3 — Diffie–Hellman Key Exchange for AES-256}
\subsection{Protocol}
Alice and Bob agree publicly on $p$ (2048 bits) and $g=2$.
\begin{enumerate}
  \item Alice chooses a random secret $a$, $1<a<p-1$, and sends $A = g^{a} \bmod p$ to Bob.
  \item Bob chooses a random secret $b$, $1<b<p-1$, and sends $B = g^{b} \bmod p$ to Alice.
  \item Alice computes $K_A = B^{a} \bmod p$, Bob computes $K_B = A^{b} \bmod p$; both equal
        $g^{ab} \bmod p$. An eavesdropper sees only $p,g,A,B$ and would need to solve the
        discrete logarithm (or the Diffie–Hellman problem) to obtain $K$.
  \item The shared secret has 2048 bits while AES-256 needs exactly 256 bits, so a key-derivation
        step is applied: $K_{AES} = \mathrm{SHA\text{-}256}(K)$, where $K$ is written as 256
        big-endian bytes. Both parties obtain the same 256-bit key.
\end{enumerate}
The exchange as such is not authenticated; in practice the values $A$ and $B$ are signed (or
certified, see the PKI laboratory) to prevent man-in-the-middle attacks.

\subsection{Values}
\noindent\textbf{Alice's secret} $a$:
\bignum{<<dh.a>>}
\noindent\textbf{Bob's secret} $b$:
\bignum{<<dh.b>>}
\noindent\textbf{Alice sends} $A = g^{a} \bmod p$:
\bignum{<<dh.A>>}
\noindent\textbf{Bob sends} $B = g^{b} \bmod p$:
\bignum{<<dh.B>>}
\noindent\textbf{Alice computes} $K_A = B^{a} \bmod p$:
\bignum{<<dh.K_alice>>}
\noindent\textbf{Bob computes} $K_B = A^{b} \bmod p$:
\bignum{<<dh.K_bob>>}
$K_A = K_B$: \textbf{<<dh.equal>>}.\\[4pt]
\noindent\textbf{AES-256 key} $K_{AES}=\mathrm{SHA\text{-}256}(K)$ (256 bits, hex):
\bignum{<<dh.aes_key_hex>>}

\subsection{Using the Key}
To demonstrate that the key is usable, Alice encrypts the message with AES-256-GCM and Bob decrypts it:
\begin{terminal}
\begin{termout}
nonce      = <<dh.aes_nonce>>
ciphertext = <<dh.aes_ciphertext>>
decrypted  = <<dh.aes_decrypted>>
\end{termout}
\end{terminal}

\section{Source Code and Repository}
\begin{center}
\url{https://github.com/don09dandidi/CS} \quad (directory \texttt{lab5/})
\end{center}
\begin{lstlisting}[style=py, caption={Core of the three tasks (excerpt from \texttt{public\_key\_crypto.py})}]
# RSA
n, phi = p * q, (p - 1) * (q - 1)
d = pow(e, -1, phi)                    # PowerMod[e, -1, phi]
c = pow(m, e, n); m2 = pow(c, d, n)    # PowerMod[m, e, n], PowerMod[c, d, n]

# ElGamal
x = random_in(2, P - 2); y = pow(G, x, P)
k = random_in(2, P - 2)
c1, c2 = pow(G, k, P), (m * pow(y, k, P)) % P
m2 = (c2 * pow(c1, P - 1 - x, P)) % P

# Diffie-Hellman
a, b = random_in(2, P - 2), random_in(2, P - 2)
A, B = pow(G, a, P), pow(G, b, P)
k_alice, k_bob = pow(B, a, P), pow(A, b, P)
aes_key = hashlib.sha256(k_alice.to_bytes(256, "big")).digest()
\end{lstlisting}

\subsection{Build and Execution Instructions}
\begin{lstlisting}[style=sh]
git clone https://github.com/don09dandidi/CS.git
cd CS/lab5
pip install cryptography          # optional, only for the AES-256-GCM demo
python3 public_key_crypto.py --save   # prints all values, saves results.json
\end{lstlisting}

\section{Conclusion}
\begin{enumerate}
  \item \textbf{RSA:} keys with a <<rsa.n_bits>>-bit modulus were generated; the message
        \texttt{"<<message>>"} was encrypted with $(e,n)$ and correctly recovered with $(d,n)$.
        Security relies on the difficulty of factoring $n$ — knowing $p,q$ immediately yields
        $\varphi(n)$ and $d$.
  \item \textbf{ElGamal:} with the 2048-bit safe prime and $g=2$ the message was encrypted into the
        pair $(c_1,c_2)$ and decrypted correctly. ElGamal is probabilistic (a fresh $k$ gives a
        different ciphertext each time) and the ciphertext is twice as long as the message.
  \item \textbf{Diffie–Hellman:} Alice and Bob obtained the same 2048-bit secret without
        transmitting it, and derived from it a 256-bit AES key with SHA-256, which was then used
        for authenticated encryption with AES-256-GCM.
  \item In practice textbook RSA/ElGamal are combined with padding schemes (OAEP) and authenticated
        key exchange; they are used to transport symmetric keys rather than to encrypt bulk data.
\end{enumerate}

\appendix
\section{Wolfram Language Code}\label{app:wl}
\lstinputlisting[style=wl]{wolfram/lab5.wl}

\begin{thebibliography}{9}
\bibitem{rsa} R.\ Rivest, A.\ Shamir, L.\ Adleman, \emph{A Method for Obtaining Digital Signatures
  and Public-Key Cryptosystems}, Communications of the ACM, 21(2), 1978.
\bibitem{elg} T.\ ElGamal, \emph{A Public Key Cryptosystem and a Signature Scheme Based on Discrete
  Logarithms}, IEEE Trans.\ Inf.\ Theory, 31(4), 1985.
\bibitem{dh} W.\ Diffie, M.\ Hellman, \emph{New Directions in Cryptography}, IEEE Trans.\ Inf.\
  Theory, 22(6), 1976.
\bibitem{rfc7919} D.\ Gillmor, \emph{RFC 7919: Negotiated Finite Field Diffie-Hellman Ephemeral
  Parameters for TLS}, 2016.
\bibitem{stallings} William Stallings, \emph{Cryptography and Network Security: Principles and
  Practice}, 7th Edition, Pearson, 2017.
\end{thebibliography}
\end{document}
